\section{Data Provenance dan Data Integrity}

Integrasi \textit{patient-generated health data} (PGHD) ke dalam sistem rekam medis elektronik terdistribusi menimbulkan kebutuhan kuat terhadap mekanisme yang mampu menelusuri asal-usul, konteks, serta evolusi data sepanjang \textit{lifecycle}-nya \parencite{ahima_fountain_2014,ahmed_dataprovenance_2023}. 
Tanpa mekanisme tersebut, penyedia layanan kesehatan akan kesulitan membedakan data yang dihasilkan langsung oleh sistem klinis dengan data yang dilaporkan pasien, keluarga, atau \textit{caregiver}, sehingga berdampak pada kepercayaan, penilaian risiko klinis, dan pengambilan keputusan berbasis data \parencite{ahima_fountain_2014,ahmed_dataprovenance_2023}. 
Pada arsitektur seperti DecMed yang mengintegrasikan PGHD melalui infrastruktur buku besar terdistribusi IOTA, isu ini menjadi semakin krusial karena data berpindah melalui banyak lapisan (perangkat IoMT, \textit{edge}, \textit{gateway}, dan layanan terpusat) sebelum akhirnya dimanfaatkan sebagai bagian dari \textit{longitudinal patient record}.

\subsection{Konsep dan Peran Data Provenance dalam PGHD}

Secara umum, data provenance didefinisikan sebagai atribut mengenai asal-usul data, proses pembentukan, serta transformasi data sepanjang siklus hidupnya \parencite{ahmed_dataprovenance_2023,jaigirdar_prov_iot_2020}. 
Dalam konteks kesehatan, provenance mencakup kemampuan organisasi untuk menelusuri sumber informasi klinis (misalnya dari \textit{health information exchange}, personal health record, atau sistem EHR), mengidentifikasi penulis dokumentasi, menentukan pemilik serta hak penggunaan data, dan melacak setiap perubahan yang terjadi sepanjang waktu \parencite{ahima_fountain_2014,ahmed_dataprovenance_2023}. 
Ahmed dkk. menekankan bahwa provenance memungkinkan pelacakan asal data, operasi yang dilakukan, dan aktor yang terlibat, sehingga menjadi basis penting untuk audit, penelusuran kesalahan, dan peningkatan kualitas data di layanan kesehatan \parencite{ahmed_dataprovenance_2023}. 

Pada PGHD, provenance menjadi kunci untuk membedakan data yang ditangkap pasien dari data yang dihasilkan sistem klinis, serta untuk menandai konteks pengumpulan seperti jenis perangkat yang digunakan, frekuensi pencatatan, atau siapa yang memasukkan data \parencite{ahima_fountain_2014,tripathy_pghd_survey_2022}. 
Fountain menegaskan bahwa kemampuan mengidentifikasi provenance PGHD yang masuk ke EHR sangat penting agar penyedia layanan dapat mengetahui data apa yang didokumentasikan oleh pasien, keluarga, atau \textit{designated caregiver}, sekaligus memutuskan bagaimana data tersebut akan diintegrasikan dalam alur kerja klinis \parencite{ahima_fountain_2014}. 
Tanpa provenance yang memadai (misalnya hanya di tingkat dokumen), granularitas sumber data yang terbatas dapat menimbulkan isu integritas dan mengurangi kepercayaan terhadap informasi yang digunakan dalam pengambilan keputusan \parencite{ahima_fountain_2014,ahmed_dataprovenance_2023}. 

\subsection{Model dan Mekanisme Data Provenance untuk IoMT dan IoT-Health}

Berbagai model provenance telah diusulkan untuk meningkatkan transparansi dan keandalan data pada lingkungan IoT dan IoMT \parencite{jaigirdar_prov_iot_2020,jaigirdar_securityaware_2023,ahmed_dataprovenance_2023}. 
W3C PROV-DM menyediakan model data standar yang memisahkan entitas data, aktivitas pemrosesan, dan agen yang bertanggung jawab, serta relasi di antaranya, sehingga cocok digunakan sebagai kerangka dasar untuk membangun graf provenance \parencite{jaigirdar_securityaware_2023}. 
Jaigirdar dkk. mengadaptasi PROV-DM untuk skenario IoT-Health dengan mendefinisikan node entitas (misalnya \textit{raw sensor data}, data teragregasi, dan hasil pemrosesan), aktivitas (registrasi perangkat, pengambilan data, propagasi, pemrosesan, dan \textit{retrieval}), serta agen (pasien, sensor, gateway, cloud, dan dokter) yang dihubungkan melalui relasi seperti \textit{wasGeneratedBy}, \textit{wasDerivedFrom}, dan \textit{wasAssociatedWith} \parencite{jaigirdar_securityaware_2023}. 

Prov-IoT kemudian memperluas model ini dengan menambahkan node \textit{security metadata} dan relasi \textit{wasProvedBy}, sehingga setiap aktivitas propagasi data memiliki bukti terkait status keamanan, konfigurasi, dan kontrol yang aktif pada saat aktivitas berlangsung \parencite{jaigirdar_prov_iot_2020,jaigirdar_securityaware_2023}. 
Pendekatan ini menghasilkan \textit{security-aware provenance graph} yang tidak hanya mendeskripsikan garis keturunan data, tetapi juga melampirkan atribut keamanan, misalnya protokol komunikasi yang digunakan, hasil verifikasi tanda tangan, status \textit{intrusion detection}, atau keberadaan \textit{web application firewall} \parencite{jaigirdar_securityaware_2023}. 
Studi kasus IoT-Health yang dikembangkan Jaigirdar dkk. menunjukkan bahwa model tersebut mampu mendeteksi skenario serangan seperti \textit{node cloning}, injeksi paket palsu, serangan \textit{denial of service}, akses tidak sah, konfigurasi WAF yang salah, hingga pelanggaran izin aplikasi dengan memvalidasi konsistensi \textit{security metadata} dibandingkan dengan kebijakan keamanan yang telah ditetapkan \parencite{jaigirdar_securityaware_2023}. 

Tinjauan sistematis oleh Ahmed dkk. mengklasifikasikan pendekatan data provenance di kesehatan ke dalam empat kategori teknologi utama: berbasis log, kriptografi, \textit{blockchain}, dan ontologi \parencite{ahmed_dataprovenance_2023}. 
Pendekatan log berfokus pada rekonstruksi riwayat operasi dari catatan kejadian; pendekatan kriptografi menekankan otentikasi dan integritas melalui tanda tangan digital dan \textit{message authentication code}; pendekatan blockchain memanfaatkan transaksi di buku besar terdistribusi untuk merekam operasi dan menyediakan jejak yang sulit dipalsukan; sedangkan pendekatan ontologi memodelkan semantik provenance menggunakan skema seperti PROV-O untuk memperkaya interoperabilitas dan analisis \parencite{ahmed_dataprovenance_2023}. 
Dalam konteks DecMed yang memanfaatkan IOTA sebagai lapisan buku besar, kombinasi graf provenance yang kaya semantik dan bukti kriptografi pada level transaksi memberikan fondasi yang kuat untuk memastikan jejak PGHD dapat diaudit dari perangkat hingga sistem EHR, tanpa bergantung pada satu titik kegagalan terpusat \parencite{ahmed_dataprovenance_2023,jaigirdar_securityaware_2023}. 

\subsection{Provenance, Transparansi, dan Trust dalam Penggunaan PGHD}

Literatur PGHD menunjukkan bahwa salah satu hambatan utama adopsi PGHD oleh klinisi adalah kurangnya kepercayaan terhadap volume, kualitas, dan keterandalan data yang dikumpulkan pasien di luar lingkungan klinis \parencite{ahima_fountain_2014,tripathy_pghd_survey_2022}. 
Fountain mencatat ketiadaan metode seragam untuk data provenance di berbagai sistem EHR dan menyoroti bahwa sebagian besar implementasi hanya mencatat provenance di level dokumen, bukan di level elemen data seperti daftar obat atau masalah klinis, sehingga menimbulkan isu integritas dan menurunkan kepercayaan terhadap informasi terintegrasi \parencite{ahima_fountain_2014}. 
Tripathy dkk. menemukan bahwa kualitas, keamanan, dan tata kelola PGHD menjadi tema dominan dalam studi integrasi PGHD ke EHR, dan menggarisbawahi perlunya mekanisme provenance yang secara eksplisit mendokumentasikan siapa yang mengumpulkan data, kapan, di mana, dan melalui perangkat apa, serta bagaimana data tersebut diproses sebelum digunakan dalam keputusan klinis \parencite{tripathy_pghd_survey_2022}. 

Penerapan \textit{security-aware provenance} pada arsitektur IoMT memberikan transparansi tambahan bagi dokter dan pihak non-teknis mengenai apakah data yang mereka lihat melewati jalur yang aman, apakah kontrol keamanan berjalan sebagaimana mestinya, dan di titik mana risiko meningkat selama propagasi data \parencite{jaigirdar_securityaware_2023}. 
Jaigirdar dkk. menunjukkan bahwa dengan menghubungkan bukti keamanan (misalnya keberhasilan verifikasi tanda tangan, status \textit{intrusion detection}, atau anomali konsumsi daya) ke setiap aktivitas dalam graf provenance, sistem dapat memberikan indikator autentikasi, integritas, ketersediaan, dan ancaman terselubung (\textit{underlying threats}) yang lebih komprehensif kepada pengguna akhir \parencite{jaigirdar_securityaware_2023}. 
Dalam konteks DecMed, pendekatan serupa dapat dimanfaatkan untuk menyajikan ringkasan validasi provenance kepada klinisi (misalnya indikator bahwa PGHD berasal dari perangkat terdaftar dan jalur komunikasi terenkripsi) sekaligus menyediakan graf lengkap bagi \textit{auditor} atau administrator keamanan untuk investigasi lebih lanjut ketika terdeteksi anomali.

\subsection{Data Integrity dalam Arsitektur Smart Healthcare}

Integritas data (\textit{data integrity}) merujuk pada jaminan bahwa data yang tersimpan atau dikomunikasikan merupakan representasi yang akurat dan tidak dimodifikasi secara tidak sah dari keadaan klinis pasien \parencite{jaigirdar_securityaware_2023, gowthamani_edgecrypto_2025}. 
Pada arsitektur smart healthcare berbasis IoMT, integritas terancam oleh berbagai serangan, antara lain injeksi paket palsu, manipulasi data selama 